\subsection{A restriction on locally free acyclic covers}

\begin{claim}{A.22}{Proved}{Locally free groups cannot supply the proposed normal-Frattini covers}
A nontrivial locally free group having a finite normal generating set
has nonzero abelianization. Consequently there is no epimorphism
$p:E\twoheadrightarrow S$ from a locally free group onto a nontrivial
perfect group $S$ with a finite normal generating set such that
\[
 N\triangleleft E,\quad N\ker(p)=E\quad\Longrightarrow\quad N=E.
 \tag{A.22.1}
\]
In particular this excludes such an epimorphism onto every nonabelian
simple group. For an arbitrary source satisfying (A.22.1) and a
nonabelian simple target $S$, every nontrivial quotient of $E$ has a
specified epimorphism onto $S$. Thus the asserted quotient-size
consequence of this proposed cover is valid, while its locally free
source is excluded by the first assertion.
\end{claim}

\begin{proof}
Suppose a nontrivial locally free group $E$ is perfect and normally
generated by $x_1,\ldots,x_r$. Delete the identity elements from this
list. For every $i$, choose a finite expression of $x_i$ as a product
of commutators $[u,v]$ in $E$. Express all the finitely many elements
$u,v$ in these expressions as finite products of conjugates of the
$x_j^{\pm1}$. Let $F$ be the subgroup generated by the $x_i$, those
$u,v$, and all the conjugators just used. It is a finitely generated
free group, by local freeness. Every chosen equation holds in this
one free group.

We give the needed separation argument explicitly. If
$t_1,\ldots,t_s$ is a free basis of $F$, the assignments
\[
 t_j\longmapsto 1+T_j
\]
define a homomorphism from $F$ to the units with constant term one
of the ring $\mathbb Z\langle\!\langle T_1,\ldots,T_s
\rangle\!\rangle$ of noncommutative formal power series. The inverse
of $1+T_j$ is its formal geometric series. This homomorphism is
injective: write a nonempty reduced word in maximal syllables as
$t_{i_1}^{e_1}\cdots t_{i_k}^{e_k}$, with $e_j\ne0$ and consecutive
$i_j$ distinct. In its expansion, the coefficient of
$T_{i_1}\cdots T_{i_k}$ is $e_1\cdots e_k\ne0$. Indeed the target
monomial has $k$ maximal variable blocks. Omitting any syllable
factor would produce at most $k-1$ such blocks; using all $k$
factors and total degree $k$ forces degree one from each factor.

Let $\mathfrak m^d$ denote the series with no term below degree
$d$. Conjugation by a unit with constant term one preserves
$1+\mathfrak m^d$, and multiplication also preserves this set.
If $U,V\in1+\mathfrak m^d$, then
$[U,V]\in1+\mathfrak m^{2d}$: modulo $\mathfrak m^{2d}$, the
positive parts of $U$ and $V$ have product zero in either order,
so the two units commute. Each $x_i$ has a nonzero positive term
by injectivity. Let $d$ be the minimum of their first positive
degrees. All selected $u,v$ have expansions in
$1+\mathfrak m^d$, because they are products of conjugates of the
$x_i^{\pm1}$. The commutator expressions therefore place every
$x_i$ in $1+\mathfrak m^{2d}$. Since $2d>d$, this contradicts
the definition of $d$. Thus $E$ cannot be perfect.

Now let $p:E\twoheadrightarrow S$ satisfy (A.22.1), with $S$
nontrivial and perfect. Since $p([E,E])=[S,S]=S$, we have
$[E,E]\ker(p)=E$, and (A.22.1) gives $[E,E]=E$. Lift a finite
normal generating set of $S$ to $E$ and let $N$ be its normal
closure. Its image is $S$, so $N\ker(p)=E$ and (A.22.1) gives
$N=E$. Hence $E$ is perfect and finitely normally generated,
contrary to the first assertion when $E$ is locally free. A
nonabelian simple group is perfect and normally generated by any
one of its nonidentity elements, proving the stated case.

Finally suppose only that $S$ is nonabelian simple and (A.22.1)
holds, without imposing local freeness on $E$. For every proper
normal subgroup $N\triangleleft E$, the normal subgroup $p(N)$
of $S$ is either trivial or all of $S$. The latter alternative
would give $N\ker(p)=E$ and hence $N=E$, a contradiction.
Thus $N\subseteq\ker(p)$, and the formula
\[
 E/N\longrightarrow S,\qquad eN\longmapsto p(e)
\]
defines the promised epimorphism. In particular a nontrivial
image of $E$ has cardinality at least $|S|$.
\end{proof}

The locally free acyclic covers in
\cite[Section~2, Proposition~2.3]{A_Berrick2010} therefore cannot
have property (A.22.1) over a nontrivial perfect finitely normally
generated target. Passing to a subgroup of such a cover does not
remove this obstruction: every finitely generated subgroup of
that subgroup is still free, so the first assertion applies again.

